\section{Discussion and related work}
\subsection{Relationship to calibration and experimental design}
\paragraph{Calibration under non-identifiability.}
Bayesian computer-model calibration distinguishes parameter uncertainty from model discrepancy \citep{kennedy2001}. Eligibility-set calibration already addresses non-identifiability in multi-agent simulators and combines distributional comparison with Bayesian optimization \citep{bai2022}. Our empirical process/trajectory gates are simpler and lack their statistical test interpretation. Bayesian optimization for likelihood-free inference allocates expensive simulator evaluations \citep{gutmann2016}; our acquisition instead allocates access to process records after rollouts have been cached. These are different resource decisions. Set calibration itself is therefore not a novelty claim. The contribution lies in its explicit coupling to stage-level evidence opportunities, composed target trajectories, and an auditable evaluation of when that coupling becomes ineffective.

\paragraph{Prediction-directed acquisition.}
Variance-based active learning has a long history \citep{cohn1996}, and Bayesian design provides a decision-theoretic framework for utility-dependent data acquisition \citep{chaloner1995}. Goal-oriented OED targets predictive quantities rather than all parameter uncertainty \citep{attia2018,zhong2026}. Our squared-covariance score is a finite binary-observation specialization of that perspective, using a variance utility rather than predictive information gain. A record with high entropy can still have zero value if every mechanism assigns it the same probability or if it is uncorrelated with the probe prediction. The relevant practical decisions are the validity of stage likelihoods, the target probe distribution, numerical support preservation, and computation reuse. The current evidence supports the last of these most clearly.

\paragraph{Social simulation systems.}
CALM combines calibrated stochastic choice with activity-based traveler simulation \citep{cheng2026}; it concerns a specific empirical mobility setting rather than the finite process-record acquisition problem studied here. SocioVerse2 develops longitudinal simulation and research loops for social systems \citep{zhang2026}. Large Population Models emphasize population-scale behavior and interactions \citep{chopra2025}. These systems motivate reusable interfaces between mechanisms, observations, and collective outcomes. They are not empirical baselines in this paper, and an adapter bearing a benchmark name is not a reproduction of their evaluation. The proposal's broader aim of recovering individual dynamics from macro observations is also reflected in macro2mind \citep{macro2mind2026}; we do not implement or compare that model. Real-domain comparisons would require compatible observation opportunities, licensed data, and a named, frozen scoring protocol.

\subsection{Scope, boundary conditions, and failure modes}
\label{sec:boundaries}
\paragraph{Concentration removes the signal used by acquisition.}
As $w$ approaches a point mass, $\Cov_w(p_r,Z)\to0$ for every record and hence $A_r(w)\to0$. A large hard-compatible set is therefore not sufficient to preserve epistemic uncertainty. The chosen small kernel width and finite-precision softmax jointly restrict effective support. Posterior updates cannot recover candidates with exactly zero mass. The observed absence of acquisition gains is consistent with this limit, but proving its quantitative contribution would require controlled interventions beyond this frozen study.

\paragraph{Why plausible development changes failed.}
Wider kernels in rounds 8--9 and 36--37 were reverted: keeping more candidate mass increased prediction error on the chosen development objective. Mixing in uniform weight in rounds 14--17 and 50--53 also failed to improve it. These outcomes distinguish numerical diversity from accurate uncertainty; they do not prove that diverse sets are generally inferior. Stronger evidence likelihoods at temperatures $0.05$ and $0.1$ worsened results in rounds 10--11, consistent with noisy binary evidence exerting excessive influence under a mismatched finite bank. Later temperature changes had only negligible effects once the mixture was already concentrated.

Rounds 18--23 added process Brier energy with almost no change in the retained behavior. Rounds 24--29 changed channel weights without exceeding the acceptance threshold. Alternative worst-channel, absolute, Huber, trimmed, and late-time objectives in rounds 31--35 did not improve the score. These losses encode different notions of fit; none is a universally valid robustness correction. Probe-coordinate rescaling in rounds 42--45 produced identical AUC, an especially direct warning against describing those rejected ideas as innovations in the final method. Seed-marginalized kernels in rounds 54--57 performed worse, but this changed both the calibration construction and kernel scale, so the result does not show that modeling randomness is harmful. Finally, raising the variance floor in rounds 59--60 reversed the benefit of round 58. Sensitivity to this regularizer is a limitation, not evidence that the selected value transfers to new domains.

\paragraph{Batch conditioning does not imply submodularity.}
Adaptive-submodularity guarantees require a diminishing-returns property \citep{golovin2011}. Our objective does not satisfy it in general. Consider a mechanism indexed by two independent fair bits $(U,V)$, target $Z=U\oplus V$, and two noiseless queries revealing $U$ and $V$. Either query alone leaves target variance unchanged, while both determine $Z$. The marginal value of the second query increases after acquiring the first, violating diminishing returns, even though query outcomes are deterministic and thus conditionally independent given the mechanism. This analytical counterexample explains why exact fantasy averaging is described as greedy conditional acquisition, without a generic approximation guarantee.

\paragraph{Agreement can coexist with arbitrary misspecification.}
Suppose all candidates predict $z_0$ while an omitted environmental rule makes reality $z_0+\Delta$. Every epistemic covariance and acquisition score is zero, while mean squared error is $\|\Delta\|^2/D$. The variance-reduction identity remains valid internally. It cannot discover a mechanism absent from the candidate class, and broad aleatoric intervals may conceal weak shifts. Post-observation residual checks provide a different source of evidence, but the tested thresholds have neither calibrated false-positive control nor advance-warning guarantees. Candidate expansion and explicit discrepancy models are possible remedies, not implemented successes.

\paragraph{What the current evidence does and does not support.}
The same development split selected 60 candidates; reporting its best score without the full search history would overstate validation. The test evaluation and independent-world extension strengthen reproducibility but retain fixed templates, training data, evidence, and a small finite mechanism bank. No confidence interval here covers uncertainty over human populations or over the search procedure itself. The Round-2 extension in Section~\ref{sec:round2} partially addresses this issue with fresh acquisition cohorts and a new sealed holdout. Further evaluation should vary event designs, historical datasets and query pools, include same-prior optimized OED, and test whether the noncollapsed acquisition benefit persists outside the constructed stress setting. It would also isolate sharing from prior changes and validate observation models before using real behavioral records.

\subsection{Conclusion}
The study makes a budgeted process-evidence loop explicit and inexpensive to execute, with complete-path prediction and separable uncertainty terms. Its strongest empirical lesson is that successful aggregate calibration can eliminate the mechanism disagreement on which active evidence selection depends. The frozen implementation improves synthetic trajectory prediction and accelerates a reference scoring calculation, but the V1 matched controls do not establish an acquisition or set-retention advantage. The separate V2 stress test restores an acquisition advantage over matched Random and Entropy while tying ParamIG; its constructed nuisance pool and inactive batch ablation limit the claim. Treating that distinction as a result yields a more useful foundation for compositional social simulation than an unsupported claim of mechanism identification or external benchmark superiority.
